\documentclass{article}
\usepackage[margin=28mm]{geometry}
\usepackage{amsmath,amssymb}
\usepackage{xurl}
\usepackage[hidelinks]{hyperref}
% Shared commands for notes in docs/paper-gaps/.

\DeclareUnicodeCharacter{2080}{\ifmmode{}_{0}\else\textsubscript{0}\fi}
\DeclareUnicodeCharacter{2081}{\ifmmode{}_{1}\else\textsubscript{1}\fi}
\DeclareUnicodeCharacter{2082}{\ifmmode{}_{2}\else\textsubscript{2}\fi}
\DeclareUnicodeCharacter{2099}{\ifmmode{}_{n}\else\textsubscript{n}\fi}

\newcommand{\Lean}{\textsc{Lean}}
\ExplSyntaxOn
\NewDocumentCommand{\leanid}{m}{
  \ifmmode
    \textsf{\detokenize{#1}}
  \else
    \group_begin:
    \urlstyle{sf}
    \tl_set:Nn \l_tmpa_tl {#1}
    \tl_replace_all:Nnn \l_tmpa_tl {\_} {_}
    \exp_args:NV \path \l_tmpa_tl
    \group_end:
  \fi
}
\ExplSyntaxOff
\providecommand{\tr}{\operatorname{tr}}
\newcommand{\tnleanIssueBase}{https://github.com/LionSR/TNLean/issues/}
\newcommand{\tnleanPrBase}{https://github.com/LionSR/TNLean/pull/}
\newcommand{\tnleanDocBase}{https://sirui-lu.com/TNLean/docs/}
\newcommand{\ghissue}[1]{\href{\tnleanIssueBase#1}{GitHub issue \##1}}
\newcommand{\ghpr}[1]{\href{\tnleanPrBase#1}{PR \##1}}
\newcommand{\doclink}[1]{\href{\tnleanDocBase#1.html}{\path{#1}}}

% Dirac notation and the identity operator, as in blueprint/src/macros/common.tex.
% \providecommand keeps note-local definitions authoritative.
\usepackage{dsfont}
\providecommand{\ket}[1]{|#1\rangle}
\providecommand{\bra}[1]{\langle #1|}
\providecommand{\braket}[2]{\langle #1 | #2 \rangle}
\providecommand{\ketbra}[2]{|#1\rangle\!\langle #2|}
\providecommand{\Id}{\mathds{1}}
\providecommand{\one}{\mathds{1}}
\providecommand{\C}{\mathbb{C}}
\providecommand{\MN}[1]{M_{#1}(\C)}
\providecommand{\MD}{M_D(\C)}

% Verdict marker: \gapnote{<kind>}{<status>}. Kind is the policy
% classification (clarification, local-correction, scope-restriction,
% unfaithful, false-source, open-gap); status is open, wip (elimination
% underway), resolved, or historical. Machine-read by the site index;
% prints a verdict line.
\newcommand{\gapnote}[2]{\par\noindent\textbf{Verdict:} #1 (#2).\par}

\providecommand{\gapnote}[2]{\par\noindent\textbf{Verdict:} #1 (#2).\par}
\providecommand{\leanid}[1]{\textsf{\detokenize{#1}}}
\title{The Padding Letter in the Power Expansion of Landau, Vazirani and Vidick}
\author{}
\date{2026-10-28}
\begin{document}
\maketitle
\gapnote{local-correction}{resolved}

\section{The source statement}
Landau, Vazirani and Vidick~\cite{Landau2015Polynomial},
\href{https://arxiv.org/abs/1307.5143v1}{arXiv:1307.5143v1}, source
\path{algo53.tex}, consider $H=\sum_{i=1}^{n-1}H_i$ on $n$ qudits, with $H_i$
acting on the qudits $i$ and $i+1$ and $0\le H_i\le 1$ (lines 283--284). They
define $A=\bigl((1-H/n)/(1-\varepsilon_0/n)\bigr)^m$ (line 600), and with
$P_i=1-H_i$ and $C=1/(1-\varepsilon_0/n)$ they write (equation
\texttt{eq:a-exp}, lines 606--609)
\begin{align*}
  A=C^m\frac{1}{n^m}\sum_{I=(i_1,\ldots,i_m)\in\{1,\ldots,n\}^m}P_I,
  \qquad P_I=\prod_{j=1}^mP_{i_j}.
\end{align*}

\section{The point at issue}
The sum runs over words in the $n$ letters $1,\ldots,n$, but $H$ has only the
$n-1$ terms $H_1,\ldots,H_{n-1}$, so $P_n$ is not defined. Expanding the power
gives $\sum_IP_I=\bigl(\sum_{i=1}^nP_i\bigr)^m$, and the identity holds exactly
when $\sum_{i=1}^nP_i=n-H$.

\section{The convention}
The expansion is read with $H_n=0$, so that $P_n=1$. Then
$\sum_{i=1}^nP_i=(n-1)-H+1=n-H$, hence
$C^mn^{-m}\sum_IP_I=\bigl(C(1-H/n)\bigr)^m=A$. The padding letter is the
identity and acts on no site, so it never crosses a cut and does not change the
rank estimate of equation \texttt{e:decomp} (lines 584--588). The same count
$n$ is the one used in the spectral estimates (lines 602--603): since
$H\le(n-1)$, the operator $1-H/n$ is positive and the denominator
$1-\varepsilon_0/n$ is positive.

The normalization divides by the number $n$ of qudits, so the statements are
made for chains with at least one site, as in the source, whose qudits are
numbered from $1$ to $n$.

\section{Formalization}
The padded bond terms and the letters are
\leanid{TNLean.MPS.AreaLaw.paddedBond} and
\leanid{TNLean.MPS.AreaLaw.bondComplement}, indexed by the $L$ sites of the
chain, with the last letter the identity. The expansion is
\leanid{TNLean.MPS.AreaLaw.powerProjector_eq_sum_bondComplementWord}, and the
rank estimate is
\leanid{TNLean.MPS.AreaLaw.operatorSchmidtRank_cutOperator_bondComplementWord_le}.
\bibliographystyle{unsrt}
\bibliography{references}
\end{document}
